\chapter{使用说明}
\begin{enumerate}
\item 首轮先读定义、定理和基础例题。合上答案后复述条件，并用一个反例说明遗漏条件的后果。
\item 第二轮按方法卡练习。每题先写“对象—条件—方法”，再计算；完成后代回、求导或量纲检查。
\item 第三轮练习多解题与竞赛题。记录哪一步利用了对称、线性、单调、正性或紧致性。
\item 章末习题每章分基础、考研风格、竞赛拓展各三题。题后答案链接与附录反向链接一一对应；先独立完成再查看。
\end{enumerate}
难度用一至三星表示相对层次。基础题用于巩固概念，考研风格题强调常见计算与综合运用，竞赛风格题强调构造、证明和推广。
\begin{method}{多解比较的记录方式}
适用条件：一道题出现两个以上自然入口。步骤：先保留各方法的关键变形，再比较运算长度、一般性和推广方向。易错点：把同一式子的重新排版当作另一种方法；不同方法应有不同的核心依据。
\end{method}
建议建立错误记录：定义域、定理条件、代数运算、边界端点、方向符号各一栏。一次订正必须回答“为什么这一步可行”，不能只记录最终答案。
